\begin{abstract}
Readable model weights are tempting evidence about the mechanisms that
generated scientific data. We audit one such candidate: the endpoint-product
score of a next-step time-series predictor. A sequence of initially plausible
architecture, benchmark, intervention, and nonlinearity stories does not
survive controls for direction conventions, lag horizons, label provenance,
data budgets, intervention targets, and classical baselines. Correcting a
source/target transpose and matching lag horizons changes Lorenz-96 PCMCI
AUROC from $0.986$ in our earlier preprint to $\CorrectedPCMCI$. On target-aware, size-matched random forcing,
all four tested methods improve, so the effect is not bottleneck-specific.
In a fixed nonlinearity confirmation, the only gate-passing bottleneck
advantage is identical at every nonlinearity strength, vanishes when the
sample size doubles, and arises because the prediction-tuned Lasso comparator
returns all-zero cross-variable scores. We demonstrate a reproducible falsification protocol on one endpoint-product
score family, with reporting and responsible-use boundaries; transfer to other
interpretability signals remains untested. This negative study illustrates how a readable
predictive representation can remain useful for hypothesis generation without
constituting causal knowledge.
\end{abstract}

\section{Introduction}

Interpretability can expose regularities that domain experts did not encode
explicitly, making it a promising interface between learned representations
and scientific discovery. But an explanation can be readable without
identifying the data-generating mechanism. In time series, predictive
association, Granger predictability, intervention response, and structural
causation answer different questions
\citep{granger1969investigating,pearl2009causality,peters2017elements}.
Confusing them can turn a useful model diagnostic into an unsupported
scientific claim.

Our motivating case is a next-step predictor with input and output
projections. For a non-lagged model, the candidate score is
\begin{equation}
  S=|W_{\mathrm{out}}W_{\mathrm{in}}|,\qquad
  S_{ij}\ \text{read as the strength of }j\rightarrow i.
  \label{eq:score}
\end{equation}
The score is cheap, compact, and often graph-like. It arose while studying a
Mamba-style state-space forecaster \citep{gu2023mamba}, but a readable pattern
does not establish that the architecture discovered causal structure. Our
earlier preprint found that a plain linear bottleneck matches or exceeds the
state-space readout \citep{lade2026prediction}, so the corrected experiments
here audit the linear factorized predictor directly.

We make three contributions. First, we organize implementation, estimand, and
provenance checks into an audit protocol for discovery claims. Second, we
apply it to corrected observational, interventional, and controlled nonlinear
experiments while retaining null and contradictory outcomes. Third, we pair
the empirical audit with a responsible-use and reproducibility boundary. We
do not propose Eq.~\ref{eq:score} as a causal-discovery method. The checks
are reusable recommendations; their empirical validation here is restricted
to one predictive score family, not interpretability methods in general.

\paragraph{Relation to prior work.}
Time-series work spans Granger/VAR methods, conditional-independence tests,
and neural estimators
\citep{runge2019pcmci,runge2020discovering,tank2021neural,nauta2019causal,pamfil2020dynotears}.
Interpretable neural Granger models explicitly design the predictor around a
causal estimand \citep{marcinkevics2021interpretable}; causal representation
learning likewise emphasizes assumptions connecting representations to
interventions \citep{scholkopf2021toward}. Our question is complementary:
what should be checked \emph{before} treating an ordinary predictive weight
as scientific evidence? Attention studies question whether readable weights
faithfully explain predictions \citep{jain2019attention}, with important
qualifications about the explanatory claim and evaluation
\citep{wiegreffe2019attention}. Saliency sanity checks show that visually
plausible maps can fail model- and data-randomization tests
\citep{adebayo2018sanity}; \citet{rudin2019stop} distinguishes inherently
interpretable models from post-hoc explanations. Our novelty is not the
warning that readability is insufficient. It is a worked audit of the
additional inference from predictive weights to causal structure, using
orientation, lag, provenance, matched-budget, and intervention-target checks.

\section{Candidate and Audit Design}

\paragraph{Lag-resolved candidate.}
Most corrected experiments use a factorized predictor of $x_t$ from
$[x_{t-1},\ldots,x_{t-p}]$. Splitting the input projection by lag gives
\[
  S_\tau=\left|W_{\mathrm{out}}W_{\mathrm{in},\tau}\right|,
  \quad \tau=1,\ldots,p,
\]
with convention $S_\tau[\mathrm{target},\mathrm{source}]$. We zero
self-edges. The predictor is linear: $W_{\mathrm{in}}$ maps the $Kp$ lagged
inputs to width $d$ and $W_{\mathrm{out}}$ maps back to $K$ outputs, with no
activation. We use $d=16$, except $d=32$ in Section~\ref{sec:nonlinearity};
because $d$ can exceed $K$, we call it factorized rather than rank-limited.
Table~\ref{tab:real} aggregates lags by their maximum score;
Section~\ref{sec:nonlinearity} instead evaluates the flattened off-diagonal
lag tensor, retaining lag identity. The ``Lagged-BN'' and ``Linear-BN'' labels (BN: bottleneck, our historical
name for this factorized predictor) refer to this prediction-trained model,
not to an intervention-aware causal model. Lasso and Ridge use the same lagged
rows and choose regularization by prediction validation rather than graph
labels.

\paragraph{Representation and protocol checks.}
We test direction, lag, and diagonal conventions on planted edges. Package
adapters are explicit: Tigramite stores values as
$[\mathrm{source},\mathrm{target},\mathrm{lag}]$, whereas our score uses
$[\mathrm{target},\mathrm{source}]$. Strictly lagged evaluation excludes
contemporaneous matrices. The original controlled comparisons match usable
observations, lag horizon, and preprocessing. The added observational baselines
share input arrays and a maximum lag horizon; VAR-LiNGAM selects its lag order
within that horizon (Appendix~\ref{app:followup}). Independent trajectories remain
separate so no lag window crosses a boundary.

\paragraph{Intervention-aware estimand.}
Each intervention segment records its target. When variable $j$ is externally
forced, the outcome equation for $j$ is excluded from that segment for every
method; other outcome equations remain usable. A combined
observational/interventional arm is compared with an observational arm having
the same number of usable lagged rows. Constant clamps, additive soft forcing,
and per-step random forcing are separate conditions.

\paragraph{Provenance and inference.}
Every graph is tagged as exact simulated, literature-derived soft, or
heuristic soft. Independent data-generating seeds are the replication unit;
model-seed spread on fixed data is only descriptive. In the original
intervention and nonlinearity studies, paired effects receive
bootstrap 95\% intervals and one-sided Wilcoxon tests with Holm correction
over each declared family. Hyperparameters and gates are fixed before graph
evaluation. The post-review trajectory comparison reports exploratory
bootstrap intervals without significance tests.

\section{Corrected Evidence}

\subsection{Benchmark provenance and orientation}

Table~\ref{tab:real} separates two observational sensitivity analyses from an
exact simulated graph. Climate contains six NOAA indices ($T=757$) and three
literature-derived teleconnection labels
\citep{trenberth1994decadal,newman2016pdo,delworth2000observed}; the
definitionally coupled ENSO/SOI pair is excluded. Finance contains ten sector
ETF return series ($T=1893$) and six heuristic labels. Lorenz-96 is a
$K=10$, $T=1500$ simulation with the exact 30-edge interaction graph
\citep{lorenz1996predictability}.

\begin{table}[t]
\centering
\small
\caption{Corrected AUROC at lag horizon 3. Climate and finance have soft
labels and are sensitivity analyses, not causal leaderboards. Lagged-BN
reports mean $\pm$ standard deviation over three model initializations on
fixed data; other rows are deterministic fits on those same data. This
is a descriptive comparison, not an estimate of across-dataset uncertainty.
Finance Lasso selects the largest penalty in its grid and returns all-zero
cross-variable scores, so its $0.500$ is a tie, not a ranking. Ridge also
selects a grid-boundary penalty on climate (smallest) and finance (largest)
but returns untied scores.}
\label{tab:real}
\begin{tabular}{lccc}
\toprule
Method & Climate & Finance & Lorenz-96 \\
\midrule
\RealBenchmarkTableRows
\bottomrule
\end{tabular}
\end{table}

Our earlier preprint \citep{lade2026prediction} reported a Lorenz-96 PCMCI
value of $0.986$ from a transposed score matrix at lag horizon 5; at the
correct orientation and the shared horizon 3 it is $\CorrectedPCMCI$. The
preprint's DYNOTEARS and VAR-LiNGAM values also included contemporaneous
coefficients (Appendix~\ref{app:followup}). Its benchmark values are
superseded by Table~\ref{tab:real} and Appendix~\ref{app:followup}.
On the exact graph, Granger and
tuned linear regression have higher observed scores than PCMCI/PCMCI+
in this configuration. The
Lagged-BN score still contains readable dependency signal, but its
$\CorrectedLaggedBNMean\pm\CorrectedLaggedBNStd$ remains below
Granger/Lasso/Ridge
($\CorrectedGranger/\CorrectedLasso/\CorrectedRidge$). Its spread over model
seeds 42--44 is initialization
stability, not independent uncertainty. Climate and finance cannot support a
causal ranking. Near-ceiling linear and Granger scores make this an
implementation and interpretation check, not a demanding test of general
causal-discovery performance. The exact reference graph describes the
continuous-time generating equations; subsampling need not preserve its
interpretation as a discrete-time conditional-independence graph.

\paragraph{Post-review trajectory replication.}
We additionally evaluate ten independently initialized Lorenz-96 trajectories
(seeds 81000--81009), retaining the original $K=10,T=1500$, forcing, sampling,
lag horizon, preprocessing, and estimator settings. Each trajectory is the
replication unit; the bottleneck averages model seeds 42--44 within it.
Table~\ref{tab:replication} reports the resulting uncertainty for all methods.
Lasso, Ridge, Granger, and VAR-LiNGAM exceed Lagged-BN with paired
intervals excluding zero; the PCMCI and PCMCI+ intervals include zero; and
DYNOTEARS is lower, with an interval excluding zero. On AUPRC, Lagged-BN is the lowest of all eight methods.
These exploratory follow-ups extend replication, not graph-family coverage.
DYNOTEARS and VAR-LiNGAM are run on the identical exported arrays in an
isolated compatible environment; neither is evaluated in the intervention
experiments. Package defaults and assumptions are detailed in
Appendix~\ref{app:followup}.

\begin{table}[t]
\centering\small
\caption{Post-review Lorenz-96 replication over ten independent trajectories.
AUROC and AUPRC are means with 95\% percentile bootstrap intervals over
trajectories; bottleneck initializations are averaged within each trajectory.
The final column is the paired AUROC difference, method minus Lagged-BN,
within each trajectory (positive favors the listed method), with its own
trajectory-bootstrap interval. All results are exploratory.}
\label{tab:replication}
\begin{tabular}{lccc}
\toprule
Method & AUROC [95\% CI] & AUPRC [95\% CI] & $\Delta$ AUROC [95\% CI] \\
\midrule
\TrajectoryTableRows
\bottomrule
\end{tabular}
\end{table}

\subsection{Intervention semantics and method specificity}

Both intervention audits use random stable linear-Gaussian VAR(3) systems
with $K\in\{10,20\}$, $T=300$ observational rows, and 10 independent seeds:
cross-edge probability $.2/3$ per ordered pair and lag, coefficients of magnitude $U(.1,.25)$ with
random sign, self-coefficients $U(.2,.4)$ at each lag, rescaled when needed
to companion spectral radius $.85$, and noise standard deviation $.1$. Each
variable is then forced for 50 steps. For method $m$, the intervention gain is
\[
 \Delta_m=\mathrm{AUROC}_m(\mathrm{combined})-
 \mathrm{AUROC}_m(\mathrm{size\mbox{-}matched\ observational}).
\]
The semantics audit tests whether $\Delta_{\mathrm{BN}}-\Delta_m>0$ for
Lasso and Ridge, using lag-resolved AUROC.
Holm-significant positive differences are only
\SemanticsPositiveRange{} AUROC. Constant clamps yield no positive
bottleneck advantage; random and soft forcing produce selected small
differences, not a general mechanism.

The four-method audit then applies target-aware random forcing to the
bottleneck, Lasso, PCMCI, and Granger, using lag-aggregated AUROC.
Figure~\ref{fig:evidence-summary} shows positive gains for every method at
both dimensions; Appendix~\ref{app:interventions} lists the values.
Combined-arm AUROC is at least $.99$ for every method, so each $\Delta_m$ is
bounded by that method's observational shortfall; between-method differences
in gain, including the Granger contrast below, partly reflect weaker
observational recovery.

\begin{figure}[t]
\centering
\includegraphics[width=\linewidth]{workshop_evidence_summary.pdf}
\caption{Corrected evidence summary. (a) Lorenz-96 AUROC against the exact
graph; Lagged-BN reports mean $\pm$ standard deviation over three model
initializations, while the baselines are deterministic. (b) Target-aware
random-forcing gain relative to a size-matched observational arm, with
bootstrap 95\% intervals over 10 independent data seeds, $T=300$.}
\label{fig:evidence-summary}
\end{figure}

All eight within-method tests reject after Holm correction
($p_{\mathrm{Holm}}=\WithinMethodHolmP$), so the gain is not
bottleneck-specific. An
exploratory paired Granger-minus-Lagged-BN contrast does not reject at $K=10$
($p_{\mathrm{Holm}}=\GrangerBNKTenP$) and is positive at $K=20$
($+\GrangerBNKTwentyDelta$,
$p_{\mathrm{Holm}}=\GrangerBNKTwentyP$). The evidence therefore supports
neither bottleneck specificity nor method invariance. Target-aware Lasso also
improves, so masking forced outcomes does not remove its gain in these
settings. Positive gains across four methods do not imply equal gains or
benefits for untested methods.

\subsection{Nonlinearity confirmation}
\label{sec:nonlinearity}

The controlled generator varies only the planted edge function
\[
 g_s(x)=a(s)\big[(1-s)x+s\sigma\tanh(x/\sigma)\big],
\]
where $\sigma\approx.085$ is the stationary source standard deviation of a
fixed Gaussian AR(1) reference and $a(s)$ equalizes marginal scale under it.
There are only two planted cross-variable edges, at lags 3 and 4;
self-dynamics, edge coefficients, noise, lag structure, and seeds remain
fixed. This is a sparse controlled check, not a broad nonlinear benchmark.
It tests our earlier preprint's exploratory finding of a bottleneck advantage
under mild nonlinearity at $K=20,T=300$ \citep{lade2026prediction}.
We compare a fixed Linear-BN ($d=32$, sparsity coefficient $10^{-3}$)
with prediction-tuned Lasso over $K\in\{10,20\}$,
$T\in\{300,600\}$, $s\in\{0,.2,.3,.4,.6,1\}$, and 10 data seeds.
The primary metric is lag-resolved AUROC at horizon 5. A cell is one
$(K,T,s)$ setting with ten paired data seeds. It passes only if all six
conditions hold: mean AUROC difference $\geq.03$; paired bootstrap 95\%
lower bound $>0$; one-sided Wilcoxon $p_{\mathrm{Holm}}<.05$ across all 24
cells; mean bottleneck AUPRC greater than Lasso's; mean bottleneck AUPRC
greater than edge prevalence; and mean held-out prediction MSE no greater
than $1.10$ times Lasso's. A regime requires the \emph{same} two or more
$s\in\{.2,.3,.4\}$ to pass at both dimensions at one common $T$.
Appendix~\ref{app:gates} gives splits, metric definitions, and all cell results.

No $T$ satisfies that rule. Six cells pass, all at $K=20,T=300$, so the
passing cells reflect one $(K,T)$ setting rather than any nonlinearity
strength. At
$K=20,T=300$, Linear-BN minus Lasso AUROC is
$+\NonlinearityKTwentyDelta$ for \emph{every} $s$, including 0 and 1. At
$K=10,T=300$, it is $+\NonlinearityKTenDelta$ for every $s$, but no cell
passes: each has bootstrap lower bound zero and adjusted $p=1$. At $T=600$,
both methods reach AUROC \NonlinearityTSixHundredAUROC{} throughout.

A post-review diagnostic explains these differences. Linear-BN attains
AUROC and AUPRC $1.00$ in all 240 fits, so every nonzero difference comes
from the comparator. Re-fitting reproduces the Lasso penalty, AUROC, and
test MSE for all 240 saved fits. At $K=20,T=300$, all 60 fits select
$\alpha=1$, the largest value in the grid, and have zero off-diagonal lag
scores; at $T=600$, all 120 fits select $\alpha=.1$ and attain perfect AUROC.
The $+\NonlinearityKTwentyDelta$ difference is thus against an uninformative
edge ranking chosen by prediction validation: comparator over-regularization
at $K=20,T=300$, not a nonlinearity-specific advantage. The design cannot
resolve a nonlinearity effect in either direction, and the gates lacked a
check that the comparator's ranking is non-degenerate; Section~\ref{sec:implications}
adds that check. No graph-label-based retuning replaces the original
comparison.

\section{Implications, Limitations, and Responsible Use}
\label{sec:implications}

\paragraph{A practical falsification checklist.}
Before elevating an interpretation to a discovery claim: (i) test direction,
lag, and diagonal conventions on planted examples; (ii) match data budgets,
preprocessing, boundaries, and seeds; (iii) encode intervention targets and
state intervention semantics; (iv) compare against sparse and classical
baselines under the same estimand, verify that each comparator gives a
non-degenerate ranking (no all-tied scores), and report any penalty selected
at a grid boundary; (v) distinguish exact graphs from soft
labels; and (vi) fix variants before evaluation, report multiplicity-adjusted
effects, and retain failures.

\paragraph{Limitations.}
The evidence covers one score family, two dimensions ($K=10,20$) in the
intervention and nonlinearity studies, one intervention sample size ($T=300$),
and two nonlinearity sample sizes ($T=300,600$). Lorenz-96 provides one
exact dynamical graph family; the controlled generator has only two planted
cross-variable edges. Climate and finance have sparse, incomplete soft
labels. Trajectory replication improves uncertainty reporting but does not
make Lorenz-96 harder or validate another interpretability signal. DYNOTEARS
and VAR-LiNGAM extend the observational comparison only; their linear-model
assumptions need not hold for Lorenz-96. Broader graph, confounding, lag,
noise, and intervention coverage remains open.

\paragraph{Responsible use and reproducibility.}
These scores must not
prescribe scientific interventions, climate policy, financial trades, or
other consequential decisions. Climate and finance labels are incomplete
hypotheses, and Lorenz-96 is simulated. Predictive dependence can guide
questions, but causal use requires domain review, explicit assumptions,
validated interventions, and evidence outside the interpreted model.
Code and result artifacts are available at
\url{https://github.com/ankitlade12/predictive-weights-audit}.
The repository's release archive contains the official template,
pinned environments, tests, commands, input hashes, submitted and post-review
result artifacts, and mechanically generated table values. New analyses are
explicitly distinguished from the original fixed-gate experiments.

\paragraph{What the checklist establishes.}
The checklist diagnoses reasons an apparent discovery can be unsupported.
Orientation fixtures test an implementation convention; they do not identify
causal structure. Matching data budgets removes a sample-count confound;
it does not equate information content. Target masking respects which
equations are externally changed; it does not guarantee intervention
identification. Classical comparisons test specificity to the interpreted
model; they do not make every baseline causally valid. Non-degeneracy checks
show that a comparison is informative; they do not make the comparator
optimal. Exact simulation
labels and soft observational hypotheses require different interpretations.
Each check targets a specific failure mode: failing one undermines a
particular discovery claim, but passing all of them does not establish
causal identification. Validating the checklist on
another interpretation signal, harder graphs, and broader confounding or
intervention settings remains future work.

\section{Conclusion}

The endpoint bottleneck is informative precisely because its causal story
looked plausible. Orientation, provenance, matched estimands, intervention
semantics, and classical baselines remove the strong interpretation. The
remaining contribution is a worked falsification protocol for one predictive
score family, with checks that can be tested on further interpretation
methods before their outputs become scientific claims.
